\section*{ЗАКЛЮЧЕНИЕ}
\addcontentsline{toc}{section}{ЗАКЛЮЧЕНИЕ}

% TODO: 2--3 страницы. Содержание:
% (а) краткое напоминание об актуальности задачи;
% (б) перечисление достигнутых результатов в нумерованном списке,
%     зеркальный задачам из введения;
% (в) перечисление защищаемых положений / новизны;
% (г) перспективы дальнейшего развития.

Локальные интернет-провайдеры в Российской Федерации испытывают потребность в недорогой автоматизации абонентского обслуживания, которая закрывала бы основной поток рутинных операций, не требуя при этом длительного внедрения и поддержки специально выделенным штатом администраторов. Существующие на рынке решения либо рассчитаны на крупного оператора, либо ушли с российского рынка, либо требуют значительной адаптации под особенности конкретного провайдера.

Для решения этой задачи в рамках выпускной квалификационной работы была разработана веб-система автоматизированного управления клиентами интернет-провайдера, интегрированная с маршрутизаторами MikroTik по протоколу RouterOS API.

Основные результаты работы:
\begin{enumerate}
\item Проведен анализ предметной области: рассмотрены ключевые бизнес-процессы интернет-провайдера, проанализированы возможности маршрутизаторов MikroTik и протокола RouterOS API, выполнен сравнительный обзор восьми существующих систем биллинга и управления абонентами, обоснована необходимость собственной разработки.
\item Сформированы функциональные и нефункциональные требования: формализованы двадцать функциональных требований, построена модель прецедентов на языке UML для шести категорий действующих лиц, сформулированы требования к надежности, информационной безопасности, а также к программной и технической совместимости.
\item Спроектирована многоуровневая клиент-серверная архитектура: выделены пять уровней ответственности (представления, обмена, программного интерфейса, бизнес-логики, хранения), описаны компоненты клиентской и серверной частей, построены диаграммы последовательности обработки запросов и развертывания.
\item Спроектирована и реализована реляционная база данных, приведенная к третьей нормальной форме: тридцать две таблицы, разделенные на шесть подсистем (управление абонентами, финансы, поддержка, доступ, инфраструктура, конфигурация), с поддержкой аудита изменений, мягкого удаления и фискальных снимков на момент выпуска документов.
\item Спроектирован и реализован программный интерфейс прикладного уровня в стиле REST: более семидесяти конечных точек, сгруппированных по ресурсам, с единым форматом ответов JSON, ролевой моделью разграничения доступа на основе разрешений в формате \texttt{<<модуль.действие>>} и сквозной журнализацией изменений.
\item Реализована подсистема интеграции с маршрутизаторами MikroTik: автоматическая синхронизация простых очередей с тарифными планами абонентов, управление правилами брандмауэра с фиксацией снимков состояния и возможностью отката, мониторинг беспроводных абонентов.
\item Реализована подсистема автоматического биллинга: фоновые задачи генерации счетов по расписанию, автоматической приостановки услуг абонентам-должникам и автоматического восстановления сервиса при погашении задолженности; диспетчер задач конфигурируется через таблицу системных настроек без вмешательства в исходный код.
\item Реализована подсистема технической поддержки абонентов: обмен сообщениями в режиме реального времени по протоколу WebSocket средствами Laravel Reverb, прикрепление файловых вложений через облачное хранилище Cloudinary, оценка качества обслуживания после закрытия обращения.
\item Проведено модульное, интеграционное и системное тестирование разработанной системы, в том числе проверка взаимодействия с реальным маршрутизатором MikroTik в стендовых условиях.
\end{enumerate}

Все требования, объявленные в техническом задании, были полностью реализованы, все задачи, поставленные в начале разработки проекта, были также решены.

Разработанная система готова к опытной эксплуатации в небольших интернет-провайдерах с абонентской базой до нескольких тысяч подключений. % TODO: добавить блок «Перспективы развития»: интеграция с FreeRADIUS,
% поддержка нескольких маршрутизаторов с балансировкой нагрузки между ними,
% мобильное приложение абонента (React Native/Flutter), модуль интеграции с
% популярными платежными агрегаторами РФ (ЮKassa, СБП), модуль рассылки
% уведомлений через Telegram-бот.

Перспективы дальнейшего развития системы:
\begin{itemize}
\item интеграция с FreeRADIUS для аутентификации абонентов по протоколам PPPoE и Hotspot;
\item поддержка распределенной инфраструктуры из нескольких маршрутизаторов MikroTik с балансировкой нагрузки и резервированием;
\item разработка мобильного приложения абонента на платформах Android и iOS;
\item подключение российских платежных агрегаторов: ЮKassa, Система быстрых платежей (СБП), <<Тинькофф Касса>>;
\item рассылка уведомлений о статусе оплаты и аварийных событиях через Telegram-бот;
\item модуль геоинформационного представления абонентской базы на интерактивной карте;
\item адаптация модуля выставления счетов под фискальные требования Российской Федерации (Федеральный закон <<О применении контрольно-кассовой техники>>).
\end{itemize}
